% exercises/common.tex —— 习题公共部分（学生版与教师版共用）
% 使用方式：student.tex / teacher.tex 先写 documentclass 并定义 ifanswers，再 input 本文件。
% 不要在这里改全书样式；字体、盒子、符号全部来自书稿的 bookstyle.sty。
\usepackage{bookstyle}
\input{pandoc-shim}
\pagestyle{fancy}
\fancyhf{}
\fancyfoot[C]{\small\thepage}
\renewcommand{\headrulewidth}{0pt}

% ---------- 答案框（仅教师版显示） ----------
\newtcolorbox{ansbox}{enhanced, breakable, colback=boxgray, colframe=boxgray,
  boxrule=0pt, leftrule=2pt, colframe=accent, arc=0pt, outer arc=0pt,
  left=8pt, right=8pt, top=4pt, bottom=4pt, before skip=4pt, after skip=8pt,
  fontupper=\small, title={答案与解析}, coltitle=ink, colbacktitle=boxgray,
  fonttitle=\small\heiti, attach title to upper, after title={\par}}

% ---------- 题目 ----------
% \problem{类型}{难度}{题干}{答案与解析}
%   类型：选择 / 填空 / 简答 / 计算 / 分析 / 设计（与书中习题的类型标签风格一致）
%   难度：基础 / 中等 / 提高
%   题号为“章号.序号”，与书中习题的编号风格一致；整份习题的章号由 \setcounter{chapter}{N} 决定
\newcounter{prob}
\newcommand{\problem}[4]{%
  \refstepcounter{prob}%
  \par\medskip\noindent{\heiti\thechapter.\arabic{prob}}\quad（#1，#2）#3\par
  \ifanswers\begin{ansbox}#4\end{ansbox}\fi}

% 选择题的四个选项：\choices{甲}{乙}{丙}{丁}，竖排，便于阅读
\newcommand{\choices}[4]{%
  \par\smallskip
  \begin{list}{}{\setlength{\leftmargin}{3.2em}\setlength{\itemsep}{1pt}\setlength{\topsep}{0pt}\setlength{\parsep}{0pt}}
    \item[(A)\hfill] #1
    \item[(B)\hfill] #2
    \item[(C)\hfill] #3
    \item[(D)\hfill] #4
  \end{list}}

% 填空线
\newcommand{\blank}{\underline{\hspace{4em}}}

% 分节小标题（题型分组用）
\newcommand{\probgroup}[1]{\par\bigskip\noindent{\heiti\large #1}\par\smallskip}

% 计算题表格：与书中例题同样的五列（步骤 / 算什么 / 算式 / 结果 / 这个数从哪来）
\newenvironment{calc}{%
  \par\smallskip\noindent\small
  \renewcommand{\arraystretch}{1.3}%
  \tabularx{\linewidth}{@{}>{\raggedright\arraybackslash}p{8mm} >{\raggedright\arraybackslash}p{17mm} Y >{\raggedright\arraybackslash}p{17mm} >{\raggedright\arraybackslash}p{36mm}@{}}
  \toprule
  \heiti 步骤 & \heiti 算什么 & \heiti 算式 & \heiti 结果 & \heiti 这个数从哪来 \\
  \midrule}{%
  \bottomrule
  \endtabularx\par\smallskip}
